% Options for packages loaded elsewhere
\PassOptionsToPackage{unicode}{hyperref}
\PassOptionsToPackage{hyphens}{url}
\documentclass[
]{article}
\usepackage{xcolor}
\usepackage[margin=1in]{geometry}
\usepackage{amsmath,amssymb}
\setcounter{secnumdepth}{-\maxdimen} % remove section numbering
\usepackage{iftex}
\ifPDFTeX
  \usepackage[T1]{fontenc}
  \usepackage[utf8]{inputenc}
  \usepackage{textcomp} % provide euro and other symbols
\else % if luatex or xetex
  \usepackage{unicode-math} % this also loads fontspec
  \defaultfontfeatures{Scale=MatchLowercase}
  \defaultfontfeatures[\rmfamily]{Ligatures=TeX,Scale=1}
\fi
\usepackage{lmodern}
\ifPDFTeX\else
  % xetex/luatex font selection
\fi
% Use upquote if available, for straight quotes in verbatim environments
\IfFileExists{upquote.sty}{\usepackage{upquote}}{}
\IfFileExists{microtype.sty}{% use microtype if available
  \usepackage[]{microtype}
  \UseMicrotypeSet[protrusion]{basicmath} % disable protrusion for tt fonts
}{}
\makeatletter
\@ifundefined{KOMAClassName}{% if non-KOMA class
  \IfFileExists{parskip.sty}{%
    \usepackage{parskip}
  }{% else
    \setlength{\parindent}{0pt}
    \setlength{\parskip}{6pt plus 2pt minus 1pt}}
}{% if KOMA class
  \KOMAoptions{parskip=half}}
\makeatother
\usepackage{color}
\usepackage{fancyvrb}
\newcommand{\VerbBar}{|}
\newcommand{\VERB}{\Verb[commandchars=\\\{\}]}
\DefineVerbatimEnvironment{Highlighting}{Verbatim}{commandchars=\\\{\}}
% Add ',fontsize=\small' for more characters per line
\usepackage{framed}
\definecolor{shadecolor}{RGB}{248,248,248}
\newenvironment{Shaded}{\begin{snugshade}}{\end{snugshade}}
\newcommand{\AlertTok}[1]{\textcolor[rgb]{0.94,0.16,0.16}{#1}}
\newcommand{\AnnotationTok}[1]{\textcolor[rgb]{0.56,0.35,0.01}{\textbf{\textit{#1}}}}
\newcommand{\AttributeTok}[1]{\textcolor[rgb]{0.13,0.29,0.53}{#1}}
\newcommand{\BaseNTok}[1]{\textcolor[rgb]{0.00,0.00,0.81}{#1}}
\newcommand{\BuiltInTok}[1]{#1}
\newcommand{\CharTok}[1]{\textcolor[rgb]{0.31,0.60,0.02}{#1}}
\newcommand{\CommentTok}[1]{\textcolor[rgb]{0.56,0.35,0.01}{\textit{#1}}}
\newcommand{\CommentVarTok}[1]{\textcolor[rgb]{0.56,0.35,0.01}{\textbf{\textit{#1}}}}
\newcommand{\ConstantTok}[1]{\textcolor[rgb]{0.56,0.35,0.01}{#1}}
\newcommand{\ControlFlowTok}[1]{\textcolor[rgb]{0.13,0.29,0.53}{\textbf{#1}}}
\newcommand{\DataTypeTok}[1]{\textcolor[rgb]{0.13,0.29,0.53}{#1}}
\newcommand{\DecValTok}[1]{\textcolor[rgb]{0.00,0.00,0.81}{#1}}
\newcommand{\DocumentationTok}[1]{\textcolor[rgb]{0.56,0.35,0.01}{\textbf{\textit{#1}}}}
\newcommand{\ErrorTok}[1]{\textcolor[rgb]{0.64,0.00,0.00}{\textbf{#1}}}
\newcommand{\ExtensionTok}[1]{#1}
\newcommand{\FloatTok}[1]{\textcolor[rgb]{0.00,0.00,0.81}{#1}}
\newcommand{\FunctionTok}[1]{\textcolor[rgb]{0.13,0.29,0.53}{\textbf{#1}}}
\newcommand{\ImportTok}[1]{#1}
\newcommand{\InformationTok}[1]{\textcolor[rgb]{0.56,0.35,0.01}{\textbf{\textit{#1}}}}
\newcommand{\KeywordTok}[1]{\textcolor[rgb]{0.13,0.29,0.53}{\textbf{#1}}}
\newcommand{\NormalTok}[1]{#1}
\newcommand{\OperatorTok}[1]{\textcolor[rgb]{0.81,0.36,0.00}{\textbf{#1}}}
\newcommand{\OtherTok}[1]{\textcolor[rgb]{0.56,0.35,0.01}{#1}}
\newcommand{\PreprocessorTok}[1]{\textcolor[rgb]{0.56,0.35,0.01}{\textit{#1}}}
\newcommand{\RegionMarkerTok}[1]{#1}
\newcommand{\SpecialCharTok}[1]{\textcolor[rgb]{0.81,0.36,0.00}{\textbf{#1}}}
\newcommand{\SpecialStringTok}[1]{\textcolor[rgb]{0.31,0.60,0.02}{#1}}
\newcommand{\StringTok}[1]{\textcolor[rgb]{0.31,0.60,0.02}{#1}}
\newcommand{\VariableTok}[1]{\textcolor[rgb]{0.00,0.00,0.00}{#1}}
\newcommand{\VerbatimStringTok}[1]{\textcolor[rgb]{0.31,0.60,0.02}{#1}}
\newcommand{\WarningTok}[1]{\textcolor[rgb]{0.56,0.35,0.01}{\textbf{\textit{#1}}}}
\usepackage{graphicx}
\makeatletter
\newsavebox\pandoc@box
\newcommand*\pandocbounded[1]{% scales image to fit in text height/width
  \sbox\pandoc@box{#1}%
  \Gscale@div\@tempa{\textheight}{\dimexpr\ht\pandoc@box+\dp\pandoc@box\relax}%
  \Gscale@div\@tempb{\linewidth}{\wd\pandoc@box}%
  \ifdim\@tempb\p@<\@tempa\p@\let\@tempa\@tempb\fi% select the smaller of both
  \ifdim\@tempa\p@<\p@\scalebox{\@tempa}{\usebox\pandoc@box}%
  \else\usebox{\pandoc@box}%
  \fi%
}
% Set default figure placement to htbp
\def\fps@figure{htbp}
\makeatother
\setlength{\emergencystretch}{3em} % prevent overfull lines
\providecommand{\tightlist}{%
  \setlength{\itemsep}{0pt}\setlength{\parskip}{0pt}}
\usepackage{bookmark}
\IfFileExists{xurl.sty}{\usepackage{xurl}}{} % add URL line breaks if available
\urlstyle{same}
\hypersetup{
  pdftitle={Kappa{,} ICC{,} Bland--Altman},
  hidelinks,
  pdfcreator={LaTeX via pandoc}}

\title{Kappa, ICC, Bland--Altman}
\author{}
\date{\vspace{-2.5em}}

\begin{document}
\maketitle

\textbf{Workflow labs} use the variant template: Goal → Approach →
Execution → Check → Report.

\subsection{Learning objectives}\label{learning-objectives}

\begin{itemize}
\tightlist
\item
  Compute Cohen's kappa for categorical agreement and explain its
  chance-correction.
\item
  Compute an intraclass correlation coefficient for continuous agreement
  and distinguish consistency from absolute agreement.
\item
  Draw a Bland--Altman plot and report limits of agreement.
\end{itemize}

\subsection{Prerequisites}\label{prerequisites}

Basic R and ggplot2.

\subsection{Background}\label{background}

Measurement-agreement studies ask whether two raters, two methods, or
two instruments give the same answer on the same units. The choice of
statistic depends on the scale of the measurement. Cohen's kappa adjusts
simple percent agreement for the agreement expected by chance given the
marginal frequencies; it ranges from −1 to 1 with common landmarks at
0.4 and 0.6. Its main weakness is sensitivity to prevalence.

For continuous measurements, the intraclass correlation (ICC) and the
Bland--Altman plot answer complementary questions. The ICC is a
single-number summary of reliability, defined in several flavours
(one-way, two-way, consistency vs absolute). The Bland--Altman plot
shows \emph{pattern}: it plots the difference between two raters against
their mean and marks the limits of agreement (typically mean ± 1.96 SD).
It reveals bias, proportional bias, and heteroscedasticity that ICCs
hide.

Reliability is not the same as agreement. Two raters can be highly
correlated (one is always twice the other) and have a terrible
agreement. Always report both and let the picture tell the pattern.

\subsection{Setup}\label{setup}

\begin{Shaded}
\begin{Highlighting}[]
\FunctionTok{library}\NormalTok{(tidyverse)}
\FunctionTok{library}\NormalTok{(broom)}
\FunctionTok{library}\NormalTok{(psych)}
\FunctionTok{library}\NormalTok{(irr)}
\FunctionTok{set.seed}\NormalTok{(}\DecValTok{42}\NormalTok{)}
\FunctionTok{theme\_set}\NormalTok{(}\FunctionTok{theme\_minimal}\NormalTok{(}\AttributeTok{base\_size =} \DecValTok{12}\NormalTok{))}
\end{Highlighting}
\end{Shaded}

\subsection{1. Goal}\label{goal}

Build two small rater datasets --- one categorical, one continuous ---
and compute the matching agreement statistics.

\subsection{2. Approach}\label{approach}

For the categorical example, simulate 100 radiograph classifications (3
categories) by two readers with substantial but not perfect agreement.
For the continuous example, simulate 60 measurements by two instruments,
one with a small constant bias.

\begin{Shaded}
\begin{Highlighting}[]
\NormalTok{cats }\OtherTok{\textless{}{-}} \FunctionTok{c}\NormalTok{(}\StringTok{"normal"}\NormalTok{, }\StringTok{"mild"}\NormalTok{, }\StringTok{"severe"}\NormalTok{)}
\NormalTok{truth }\OtherTok{\textless{}{-}} \FunctionTok{sample}\NormalTok{(cats, }\DecValTok{100}\NormalTok{, }\AttributeTok{replace =} \ConstantTok{TRUE}\NormalTok{, }\AttributeTok{prob =} \FunctionTok{c}\NormalTok{(}\FloatTok{0.5}\NormalTok{, }\FloatTok{0.3}\NormalTok{, }\FloatTok{0.2}\NormalTok{))}
\NormalTok{r1 }\OtherTok{\textless{}{-}} \FunctionTok{ifelse}\NormalTok{(}\FunctionTok{runif}\NormalTok{(}\DecValTok{100}\NormalTok{) }\SpecialCharTok{\textless{}} \FloatTok{0.2}\NormalTok{, }\FunctionTok{sample}\NormalTok{(cats, }\DecValTok{100}\NormalTok{, }\AttributeTok{replace =} \ConstantTok{TRUE}\NormalTok{), truth)}
\NormalTok{r2 }\OtherTok{\textless{}{-}} \FunctionTok{ifelse}\NormalTok{(}\FunctionTok{runif}\NormalTok{(}\DecValTok{100}\NormalTok{) }\SpecialCharTok{\textless{}} \FloatTok{0.25}\NormalTok{, }\FunctionTok{sample}\NormalTok{(cats, }\DecValTok{100}\NormalTok{, }\AttributeTok{replace =} \ConstantTok{TRUE}\NormalTok{), truth)}
\NormalTok{kap\_tbl }\OtherTok{\textless{}{-}} \FunctionTok{tibble}\NormalTok{(}\AttributeTok{r1 =} \FunctionTok{factor}\NormalTok{(r1, }\AttributeTok{levels =}\NormalTok{ cats),}
                  \AttributeTok{r2 =} \FunctionTok{factor}\NormalTok{(r2, }\AttributeTok{levels =}\NormalTok{ cats))}

\CommentTok{\# continuous}
\NormalTok{n }\OtherTok{\textless{}{-}} \DecValTok{60}
\NormalTok{true\_val }\OtherTok{\textless{}{-}} \FunctionTok{rnorm}\NormalTok{(n, }\DecValTok{100}\NormalTok{, }\DecValTok{15}\NormalTok{)}
\NormalTok{inst1 }\OtherTok{\textless{}{-}}\NormalTok{ true\_val }\SpecialCharTok{+} \FunctionTok{rnorm}\NormalTok{(n, }\DecValTok{0}\NormalTok{, }\DecValTok{3}\NormalTok{)}
\NormalTok{inst2 }\OtherTok{\textless{}{-}}\NormalTok{ true\_val }\SpecialCharTok{+} \DecValTok{2} \SpecialCharTok{+} \FunctionTok{rnorm}\NormalTok{(n, }\DecValTok{0}\NormalTok{, }\DecValTok{3}\NormalTok{)    }\CommentTok{\# small positive bias}
\NormalTok{meas }\OtherTok{\textless{}{-}} \FunctionTok{tibble}\NormalTok{(inst1, inst2)}
\end{Highlighting}
\end{Shaded}

\subsection{3. Execution}\label{execution}

Cohen's kappa:

\begin{Shaded}
\begin{Highlighting}[]

\end{Highlighting}
\end{Shaded}

ICC via \texttt{psych}:

\begin{Shaded}
\begin{Highlighting}[]

\end{Highlighting}
\end{Shaded}

Bland--Altman:

\begin{Shaded}
\begin{Highlighting}[]
  \FunctionTok{mutate}\NormalTok{(}\AttributeTok{mean\_val =}\NormalTok{ (inst1 }\SpecialCharTok{+}\NormalTok{ inst2) }\SpecialCharTok{/} \DecValTok{2}\NormalTok{,}
         \AttributeTok{diff\_val =}\NormalTok{ inst2 }\SpecialCharTok{{-}}\NormalTok{ inst1)}
\NormalTok{loa }\OtherTok{\textless{}{-}} \FunctionTok{mean}\NormalTok{(ba}\SpecialCharTok{$}\NormalTok{diff\_val) }\SpecialCharTok{+} \FunctionTok{c}\NormalTok{(}\SpecialCharTok{{-}}\FloatTok{1.96}\NormalTok{, }\DecValTok{0}\NormalTok{, }\FloatTok{1.96}\NormalTok{) }\SpecialCharTok{*} \FunctionTok{sd}\NormalTok{(ba}\SpecialCharTok{$}\NormalTok{diff\_val)}

\FunctionTok{ggplot}\NormalTok{(ba, }\FunctionTok{aes}\NormalTok{(mean\_val, diff\_val)) }\SpecialCharTok{+}
  \FunctionTok{geom\_point}\NormalTok{(}\AttributeTok{alpha =} \FloatTok{0.7}\NormalTok{) }\SpecialCharTok{+}
  \FunctionTok{geom\_hline}\NormalTok{(}\AttributeTok{yintercept =}\NormalTok{ loa[}\DecValTok{1}\NormalTok{], }\AttributeTok{linetype =} \DecValTok{2}\NormalTok{, }\AttributeTok{colour =} \StringTok{"firebrick"}\NormalTok{) }\SpecialCharTok{+}
  \FunctionTok{geom\_hline}\NormalTok{(}\AttributeTok{yintercept =}\NormalTok{ loa[}\DecValTok{2}\NormalTok{], }\AttributeTok{linetype =} \DecValTok{1}\NormalTok{, }\AttributeTok{colour =} \StringTok{"steelblue"}\NormalTok{) }\SpecialCharTok{+}
  \FunctionTok{geom\_hline}\NormalTok{(}\AttributeTok{yintercept =}\NormalTok{ loa[}\DecValTok{3}\NormalTok{], }\AttributeTok{linetype =} \DecValTok{2}\NormalTok{, }\AttributeTok{colour =} \StringTok{"firebrick"}\NormalTok{) }\SpecialCharTok{+}
  \FunctionTok{labs}\NormalTok{(}\AttributeTok{x =} \StringTok{"Mean of two instruments"}\NormalTok{,}
       \AttributeTok{y =} \StringTok{"Difference (inst2 − inst1)"}\NormalTok{)}
\end{Highlighting}
\end{Shaded}

\subsection{4. Check}\label{check}

The ICC should be high (\textgreater{} 0.9) because the raters are well
correlated, but the Bland--Altman plot shows a small positive bias
(inst2 reads about 2 units higher on average).

\subsection{5. Report}\label{report}

\begin{quote}
Cohen's kappa for the two radiograph readers was
\texttt{round(kappa2(kap\_tbl{[},\ c("r1","r2"){]})\$value,\ 2)}. For
the two instruments, the ICC (absolute agreement, two-way random) was
\texttt{round(ICC(as.matrix(meas))\$results\$ICC{[}2{]},\ 2)}, but the
Bland--Altman plot revealed a mean bias of
\texttt{round(mean(ba\$diff\_val),\ 1)} units with 95\% limits of
agreement from \texttt{round(loa{[}1{]},\ 1)} to
\texttt{round(loa{[}3{]},\ 1)}.
\end{quote}

\subsection{Common pitfalls}\label{common-pitfalls}

\begin{itemize}
\tightlist
\item
  Reporting percent agreement instead of kappa.
\item
  Using Pearson \emph{r} on two raters and calling it agreement.
\item
  Omitting the limits of agreement from a Bland--Altman plot.
\end{itemize}

\subsection{Further reading}\label{further-reading}

\begin{itemize}
\tightlist
\item
  Bland JM, Altman DG (1986), \emph{Statistical methods for assessing
  agreement\ldots{}}
\item
  Shrout PE, Fleiss JL (1979), \emph{Intraclass correlations\ldots{}}
\item
  McGraw KO, Wong SP (1996), \emph{Forming inferences about some ICCs}.
\end{itemize}

\subsection{Session info}\label{session-info}

\begin{Shaded}
\begin{Highlighting}[]

\end{Highlighting}
\end{Shaded}

\subsection{See also --- chapter index}\label{see-also-chapter-index}

\begin{itemize}
\tightlist
\item
  \href{../index.Rmd}{Methods in this chapter}
\item
  \href{../../../ch00-find-your-method.Rmd}{Find your method (Chapter
  0)}
\end{itemize}

\end{document}
